% Generated draft: requires mathematical and source review.
\documentclass[11pt]{article}
\usepackage[margin=25mm]{geometry}
\usepackage{fontspec}
\setmainfont{Latin Modern Roman}
\usepackage{amsmath,amssymb,mathtools,unicode-math}
\setmathfont{Latin Modern Math}
\usepackage{enumitem,needspace,xurl,hyperref}
\setlength{\parindent}{0pt}
\setlength{\parskip}{5pt}
\setlength{\emergencystretch}{3em}
\newcommand{\N}{\mathbb{N}}
\newcommand{\Z}{\mathbb{Z}}
\newcommand{\Q}{\mathbb{Q}}
\newcommand{\R}{\mathbb{R}}
\newcommand{\C}{\mathbb{C}}
\newcommand{\F}{\mathbb{F}}
\newcommand{\A}{\mathbb{A}}
\newcommand{\PP}{\mathbb{P}}
\newcommand{\E}{\mathbb{E}}
\newcommand{\eps}{\varepsilon}
\newcommand{\dd}{\,\mathrm{d}}
\begin{document}
\section*{UnsolvedMath ID: 30001829}
\textbf{Source problem code:} catalogue entry
\subsection{Definitions and mathematical statement}
The supplied question asks: For a mechanical system driven by one or more thermostats, does an invariant measure exist, is it unique, and do reasonable initial distributions converge to it, with what rates?

The record does not specify the mechanical system, thermostat equations, phase space, evolution, admissible measures, initial distributions, or meaning of convergence. Consequently, a precise system-specific mathematical statement cannot be recovered. The following gives only a conditional probability-measure formulation of the questions; it does not supply the missing hypotheses.

Suppose a specified thermostated system has a measurable state space \((X,\Sigma)\) and a well-defined evolution \(U_t\) of probability laws on \(X\). Write \(\mathcal P(X,\Sigma)\) for the probability measures on this space. The time parameter \(t\) must range over the time domain of the specified system; neither continuous nor discrete time is selected here. A probability measure \(\mu\) is invariant if
\[
U_t\mu=\mu\qquad\text{for every admissible time }t.
\]
Once a class \(\mathcal M\subseteq\mathcal P(X,\Sigma)\) of candidate invariant measures is specified, define
\[
\mathcal I=\{\mu\in\mathcal M: U_t\mu=\mu\text{ for every admissible }t\}.
\]
The existence and uniqueness questions then ask, respectively, whether \(\mathcal I\neq\varnothing\) and whether \(\mathcal I\) contains exactly one measure. If there is a unique such measure \(\mu_*\), the convergence question requires a specified class \(\mathcal A\subseteq\mathcal P(X,\Sigma)\) of initial laws and a specified mode \(\tau\) of convergence, and asks whether
\[
U_t\nu\longrightarrow\mu_*\quad\text{in the mode }\tau
\qquad(t\to\infty),\qquad\text{for every }\nu\in\mathcal A.
\]
The rate question asks for quantitative estimates of this convergence in a specified discrepancy or against specified test functions, with the dependence on the initial law and system parameters stated explicitly. The classes \(\mathcal M\) and \(\mathcal A\), the mode \(\tau\), and any quantitative discrepancy remain unspecified. The source also does not expressly require the invariant measure itself to be a probability measure, so that normalization is conditional in this template.

There is a material source qualification: the supplied status and dated literature assessment label the problem open, whereas the supplied statement-verification note says that the official report gives affirmative results for existence, uniqueness, convergence, and rates for the thermostated system under discussion. The record labels the statement \texttt{unrecoverable} and provides no clean replacement statement. Neither the applicable affirmative theorem nor a residual open problem can therefore be stated precisely from this record.
\subsection{Short English statement}
The record asks whether thermostated mechanical dynamics have an invariant measure, whether it is unique, and whether suitable initial laws approach it at quantifiable rates. The system and essential hypotheses are missing. Its supplied open label also conflicts with a verification note reporting affirmative results in the cited report; the intended unresolved question requires human review.
\subsection{Sources}
\begin{itemize}
\item Supplied local problem registry record, permanent UnsolvedMath ID \(30001829\), source code \texttt{catalogue entry}, title \emph{Invariant Measures for Thermostated Mechanical Systems}, in the Oberwolfach Reports Open Problems collection. The record gives proposed year 2011 and names no proposer. Its recorded update timestamp is \texttt{2026-09-16T12:06:53Z}.
\item Supplied source citation: \emph{Dynamische Systeme (2012)}, Oberwolfach Reports, DOI \texttt{10.4171/owr/2011/34}, pp.\ 1913--1970. Supplied source URL and literature-source link, labeled ``original Oberwolfach report'': \url{https://doi.org/10.4171/owr/2011/34}. The supplied background describes Oberwolfach Reports as documenting mathematical workshops and their open-problem discussions. No particular theorem, contribution, or problem page is identified.
\item Literature assessment dated exactly \texttt{2026-08-21}: supplied dated curation, not independently verified. It labels the problem open and says that the extracted invariant-measure and convergence questions lack a specified dynamical system and that no general answer is available from the cited source. The background expressly calls for independent verification of time-sensitive status and cited claims.
\item Supplied statement-verification qualification: the official report is said to give affirmative results for invariant-measure existence, uniqueness, convergence, and rates in the thermostated system under discussion. The record retains the original publicly visible statement for provenance, promotes no clean statement, and marks its statement status \texttt{unrecoverable}. This qualification is supplied metadata, not an independently verified result.
\end{itemize}
\end{document}
